\textbf{Lemma.} \uline{Let} @@M0@@ \uline{be a projective and perfect morphism}, \uline{with} @@M1@@ \uline{quasi-compact}, \uline{and let} @@M2@@ \uline{be a perfect complex over} @@M3@@, \uline{acyclic in degrees} @@M4@@. \uline{Then there exists an integer} @@M5@@ \uline{such that for} @@M6@@ \uline{the complex} @@M7@@ \uline{is acyclic in degrees} @@M8@@.

Apply this lemma to the complex reduced to @@M9@@; there is therefore an integer @@M10@@ such that for @@M11@@, we have @@M12@@ when @@M13@@. We then descend by induction to @@M14@@. For this, consider the exact sequence:
@@M15@@
and apply the functor @@M16@@ to it. We have @@M17@@. Since @@M18@@, @@M19@@ is @@M20@@-acyclic for @@M21@@. Thus for every @@M22@@, and every @@M23@@:
@@M24@@
and consequently we have the exact sequences:
@@M25@@
which complete the induction.

ii) We must show that the canonical morphism:
@@M26@@
is an isomorphism. Since the property is local, we may suppose that @@M27@@ is defined by a sequence @@M28@@ of sections of @@M29@@ such that the corresponding homomorphism @@M30@@ is regular. Put: @@M31@@; we define a surjective algebra homomorphism: @@M32@@, by sending
